% Generado por generar_datos.py: no editar a mano.
\begin{minipage}[t]{0.32\textwidth}\centering
\begin{tikzpicture}[x=0.92cm,y=0.92cm]
\coordinate (c1) at (0,0);
\coordinate (c2) at (1,0);
\coordinate (c3) at (2,0);
\coordinate (c4) at (0,1);
\coordinate (c5) at (1,1);
\coordinate (c6) at (2,1);
\coordinate (c7) at (0,2);
\coordinate (c8) at (1,2);
\coordinate (c9) at (2,2);
\draw[arista] (c1) -- (c2);
\draw[arista] (c1) -- (c4);
\draw[arista] (c1) -- (c5);
\draw[arista] (c2) -- (c3);
\draw[arista] (c2) -- (c4);
\draw[arista] (c2) -- (c5);
\draw[arista] (c2) -- (c6);
\draw[arista] (c3) -- (c5);
\draw[arista] (c3) -- (c6);
\draw[arista] (c4) -- (c5);
\draw[arista] (c4) -- (c7);
\draw[arista] (c4) -- (c8);
\draw[arista] (c5) -- (c6);
\draw[arista] (c5) -- (c7);
\draw[arista] (c5) -- (c8);
\draw[arista] (c5) -- (c9);
\draw[arista] (c6) -- (c8);
\draw[arista] (c6) -- (c9);
\draw[arista] (c7) -- (c8);
\draw[arista] (c8) -- (c9);
\node[nodosig] at (c1) {1};
\node[grado] at (0.33,0.3) {3};
\node[nodo] at (c2) {2};
\node[grado] at (1.33,0.3) {5};
\node[nodo] at (c3) {3};
\node[grado] at (2.33,0.3) {3};
\node[nodo] at (c4) {4};
\node[grado] at (0.33,1.3) {5};
\node[nodo] at (c5) {5};
\node[grado] at (1.33,1.3) {8};
\node[nodo] at (c6) {6};
\node[grado] at (2.33,1.3) {5};
\node[nodo] at (c7) {7};
\node[grado] at (0.33,2.3) {3};
\node[nodo] at (c8) {8};
\node[grado] at (1.33,2.3) {5};
\node[nodo] at (c9) {9};
\node[grado] at (2.33,2.3) {3};
\end{tikzpicture}
\par\smallskip{\footnotesize (a) Inicio: esquinas 3, lados 5, centro 8. Se elige 1.}\end{minipage}\hfill
\begin{minipage}[t]{0.32\textwidth}\centering
\begin{tikzpicture}[x=0.92cm,y=0.92cm]
\coordinate (c1) at (0,0);
\coordinate (c2) at (1,0);
\coordinate (c3) at (2,0);
\coordinate (c4) at (0,1);
\coordinate (c5) at (1,1);
\coordinate (c6) at (2,1);
\coordinate (c7) at (0,2);
\coordinate (c8) at (1,2);
\coordinate (c9) at (2,2);
\draw[arista] (c2) -- (c3);
\draw[arista] (c2) -- (c4);
\draw[arista] (c2) -- (c5);
\draw[arista] (c2) -- (c6);
\draw[arista] (c3) -- (c5);
\draw[arista] (c3) -- (c6);
\draw[arista] (c4) -- (c5);
\draw[arista] (c4) -- (c7);
\draw[arista] (c4) -- (c8);
\draw[arista] (c5) -- (c6);
\draw[arista] (c5) -- (c7);
\draw[arista] (c5) -- (c8);
\draw[arista] (c5) -- (c9);
\draw[arista] (c6) -- (c8);
\draw[arista] (c6) -- (c9);
\draw[arista] (c7) -- (c8);
\draw[arista] (c8) -- (c9);
\node[nodoelim] at (c1) {1};
\node[nodo] at (c2) {2};
\node[grado] at (1.33,0.3) {4};
\node[nodosig] at (c3) {3};
\node[grado] at (2.33,0.3) {3};
\node[nodo] at (c4) {4};
\node[grado] at (0.33,1.3) {4};
\node[nodo] at (c5) {5};
\node[grado] at (1.33,1.3) {7};
\node[nodo] at (c6) {6};
\node[grado] at (2.33,1.3) {5};
\node[nodo] at (c7) {7};
\node[grado] at (0.33,2.3) {3};
\node[nodo] at (c8) {8};
\node[grado] at (1.33,2.3) {5};
\node[nodo] at (c9) {9};
\node[grado] at (2.33,2.3) {3};
\end{tikzpicture}
\par\smallskip{\footnotesize (b) Sin 1 (sin relleno). Se elige 3.}\end{minipage}\hfill
\begin{minipage}[t]{0.32\textwidth}\centering
\begin{tikzpicture}[x=0.92cm,y=0.92cm]
\coordinate (c1) at (0,0);
\coordinate (c2) at (1,0);
\coordinate (c3) at (2,0);
\coordinate (c4) at (0,1);
\coordinate (c5) at (1,1);
\coordinate (c6) at (2,1);
\coordinate (c7) at (0,2);
\coordinate (c8) at (1,2);
\coordinate (c9) at (2,2);
\draw[arista] (c2) -- (c4);
\draw[arista] (c2) -- (c5);
\draw[arista] (c2) -- (c6);
\draw[arista] (c4) -- (c5);
\draw[arista] (c4) -- (c7);
\draw[arista] (c4) -- (c8);
\draw[arista] (c5) -- (c6);
\draw[arista] (c5) -- (c7);
\draw[arista] (c5) -- (c8);
\draw[arista] (c5) -- (c9);
\draw[arista] (c6) -- (c8);
\draw[arista] (c6) -- (c9);
\draw[arista] (c7) -- (c8);
\draw[arista] (c8) -- (c9);
\node[nodoelim] at (c1) {1};
\node[nodosig] at (c2) {2};
\node[grado] at (1.33,0.3) {3};
\node[nodoelim] at (c3) {3};
\node[nodo] at (c4) {4};
\node[grado] at (0.33,1.3) {4};
\node[nodo] at (c5) {5};
\node[grado] at (1.33,1.3) {6};
\node[nodo] at (c6) {6};
\node[grado] at (2.33,1.3) {4};
\node[nodo] at (c7) {7};
\node[grado] at (0.33,2.3) {3};
\node[nodo] at (c8) {8};
\node[grado] at (1.33,2.3) {5};
\node[nodo] at (c9) {9};
\node[grado] at (2.33,2.3) {3};
\end{tikzpicture}
\par\smallskip{\footnotesize (c) Sin 3 (sin relleno). Empate 2, 7, 9: se elige 2.}\end{minipage}

\bigskip
\begin{minipage}[t]{0.32\textwidth}\centering
\begin{tikzpicture}[x=0.92cm,y=0.92cm]
\coordinate (c1) at (0,0);
\coordinate (c2) at (1,0);
\coordinate (c3) at (2,0);
\coordinate (c4) at (0,1);
\coordinate (c5) at (1,1);
\coordinate (c6) at (2,1);
\coordinate (c7) at (0,2);
\coordinate (c8) at (1,2);
\coordinate (c9) at (2,2);
\draw[arista] (c4) -- (c5);
\draw[relleno] (c4) to[bend left=30] (c6);
\draw[arista] (c4) -- (c7);
\draw[arista] (c4) -- (c8);
\draw[arista] (c5) -- (c6);
\draw[arista] (c5) -- (c7);
\draw[arista] (c5) -- (c8);
\draw[arista] (c5) -- (c9);
\draw[arista] (c6) -- (c8);
\draw[arista] (c6) -- (c9);
\draw[arista] (c7) -- (c8);
\draw[arista] (c8) -- (c9);
\node[nodoelim] at (c1) {1};
\node[nodoelim] at (c2) {2};
\node[nodoelim] at (c3) {3};
\node[nodo] at (c4) {4};
\node[grado] at (0.33,1.3) {4};
\node[nodo] at (c5) {5};
\node[grado] at (1.33,1.3) {5};
\node[nodo] at (c6) {6};
\node[grado] at (2.33,1.3) {4};
\node[nodosig] at (c7) {7};
\node[grado] at (0.33,2.3) {3};
\node[nodo] at (c8) {8};
\node[grado] at (1.33,2.3) {5};
\node[nodo] at (c9) {9};
\node[grado] at (2.33,2.3) {3};
\end{tikzpicture}
\par\smallskip{\footnotesize (d) Sin 2: aparece la arista 4--6. Se elige 7.}\end{minipage}\hfill
\begin{minipage}[t]{0.32\textwidth}\centering
\begin{tikzpicture}[x=0.92cm,y=0.92cm]
\coordinate (c1) at (0,0);
\coordinate (c2) at (1,0);
\coordinate (c3) at (2,0);
\coordinate (c4) at (0,1);
\coordinate (c5) at (1,1);
\coordinate (c6) at (2,1);
\coordinate (c7) at (0,2);
\coordinate (c8) at (1,2);
\coordinate (c9) at (2,2);
\draw[arista] (c4) -- (c5);
\draw[relleno] (c4) to[bend left=30] (c6);
\draw[arista] (c4) -- (c8);
\draw[arista] (c5) -- (c6);
\draw[arista] (c5) -- (c8);
\draw[arista] (c5) -- (c9);
\draw[arista] (c6) -- (c8);
\draw[arista] (c6) -- (c9);
\draw[arista] (c8) -- (c9);
\node[nodoelim] at (c1) {1};
\node[nodoelim] at (c2) {2};
\node[nodoelim] at (c3) {3};
\node[nodosig] at (c4) {4};
\node[grado] at (0.33,1.3) {3};
\node[nodo] at (c5) {5};
\node[grado] at (1.33,1.3) {4};
\node[nodo] at (c6) {6};
\node[grado] at (2.33,1.3) {4};
\node[nodoelim] at (c7) {7};
\node[nodo] at (c8) {8};
\node[grado] at (1.33,2.3) {4};
\node[nodo] at (c9) {9};
\node[grado] at (2.33,2.3) {3};
\end{tikzpicture}
\par\smallskip{\footnotesize (e) Sin 7 (sin relleno). Se elige 4.}\end{minipage}\hfill
\begin{minipage}[t]{0.32\textwidth}\centering
\begin{tikzpicture}[x=0.92cm,y=0.92cm]
\coordinate (c1) at (0,0);
\coordinate (c2) at (1,0);
\coordinate (c3) at (2,0);
\coordinate (c4) at (0,1);
\coordinate (c5) at (1,1);
\coordinate (c6) at (2,1);
\coordinate (c7) at (0,2);
\coordinate (c8) at (1,2);
\coordinate (c9) at (2,2);
\draw[arista] (c5) -- (c6);
\draw[arista] (c5) -- (c8);
\draw[arista] (c5) -- (c9);
\draw[arista] (c6) -- (c8);
\draw[arista] (c6) -- (c9);
\draw[arista] (c8) -- (c9);
\node[nodoelim] at (c1) {1};
\node[nodoelim] at (c2) {2};
\node[nodoelim] at (c3) {3};
\node[nodoelim] at (c4) {4};
\node[nodosig] at (c5) {5};
\node[grado] at (1.33,1.3) {3};
\node[nodo] at (c6) {6};
\node[grado] at (2.33,1.3) {3};
\node[nodoelim] at (c7) {7};
\node[nodo] at (c8) {8};
\node[grado] at (1.33,2.3) {3};
\node[nodo] at (c9) {9};
\node[grado] at (2.33,2.3) {3};
\end{tikzpicture}
\par\smallskip{\footnotesize (f) Sin 4: 5, 6, 8 y 9 ya forman un bloque lleno.}\end{minipage}
